\documentclass[a4paper,11pt]{article}
\usepackage[utf8]{inputenc}
\usepackage[T1]{fontenc}
\usepackage[french]{babel}
\usepackage[left=2cm,right=2cm,top=2cm,bottom=2cm]{geometry}
\usepackage{xcolor}
\usepackage{hyperref}
\usepackage{fontawesome5}
\usepackage{parskip}

% --- Couleurs La Banque Postale ---
\definecolor{primary}{RGB}{0, 82, 147}
\hypersetup{colorlinks=true, urlcolor=primary}

\begin{document}

% --- EXPÉDITEUR ---
\noindent
\begin{minipage}[t]{0.5\textwidth}
    \textbf{Loïc Aron MBASSI EWOLO}\\
    Élève-Ingénieur 2A -- Double diplôme ENSEA\\[3pt]
    \faMapMarker\ Cergy, France\\
    \faPhone\ (+33) 7 59 52 33 98\\
    \faEnvelope\ \href{mailto:loicaron.mbassiewolo@ensea.fr}{loicaron.mbassiewolo@ensea.fr}
\end{minipage}
\hfill
% --- DATE ---
\begin{minipage}[t]{0.4\textwidth}
    \raggedleft
    Cergy, le 6 février 2026
\end{minipage}

\vspace{1.2cm}

% --- DESTINATAIRE ---
\hfill
\begin{minipage}[t]{0.5\textwidth}
    \raggedleft
    \textbf{DSI de La Banque Postale}\\
    Direction Architecture \& Innovation\\
    Gradignan, France
\end{minipage}

\vspace{1cm}

\noindent\textbf{Objet :} Candidature au stage Développeuse/Développeur JAVA (4 à 6 mois)

\vspace{0.8cm}

Madame, Monsieur,

Votre POC d'intégration d'IA générative dans une stack Java Spring Boot / Angular rejoint un projet que j'ai mené récemment : un système multi-agents avec API LLM, RAG et déploiement Docker. Cette expérience m'a donné envie de rejoindre la Squad Environnement de développement de la DSIBA.

J'ai conçu un système multi-agents intégrant une API LLM (Gemini 2.0) avec enrichissement contextuel par RAG. Ce projet m'a conduit à développer une API REST pour l'interfaçage avec le moteur d'IA, à gérer la contextualisation des échanges utilisateur/IA, et à déployer l'ensemble dans un environnement Docker. Ces compétences sont transférables aux objectifs de votre stage : étude du module IA de Spring Boot, interfaçage avec votre API Gateway, et déploiement sur socle Cloud Privé.

Mon parcours en développement Java (Design Patterns, architecture MVC, certification POO EPFL) et en frameworks front-end SPA (React/Redux, composants réutilisables) me permettra de contribuer à l'IHM Angular. J'ai également travaillé 4 mois chez CSC, une fintech bancaire, où j'ai développé une plateforme de transactions sécurisée avec gestion des accès concurrents et conformité réglementaire.

Je suis habitué au travail en équipe Agile : sprints, revues de code, et itérations régulières. L'organisation Scrum de votre tribu Moyen de Développement correspond à ma méthode de travail. Je suis disponible dès avril 2026 pour 4 à 6 mois et mobile sur Gradignan.

Je reste à votre disposition pour un entretien afin de discuter de ma candidature et de ma contribution potentielle à votre projet de Développeur Augmenté.

Je vous prie d'agréer, Madame, Monsieur, l'expression de mes salutations distinguées.

\vspace{0.8cm}

\textbf{Loïc Aron MBASSI EWOLO}\\
\textit{Élève-Ingénieur ENSEA / Polytechnique Yaoundé}

\end{document}
